\subsection{Methodology}
\label{sec:method}

This chapter defines the experimental methodology used to answer the research questions.
The emphasis is on controlled comparison rather than on maximising the performance of any one attack variant in isolation.
To make the results interpretable, the study fixes the dataset, classifier, eligibility rule, and reporting pipeline, and then varies the attack procedure while holding the rest of the evaluation framework constant.

\subsection{Methodological overview}
The methodology is designed around three linked evaluation axes.
The first is \emph{attack effectiveness}: whether a one-pixel perturbation can overturn a correct clean prediction.
The second is \emph{attack cost}: how many model queries are required to obtain that outcome.
The third is \emph{explanation behaviour}: how explanation maps change between clean and adversarial inputs when the attack succeeds.

This three-way framing matters because a guidance strategy could appear useful under one axis while failing under another.
For example, a guided attack could produce more stable explanations on its successful cases while simultaneously lowering the overall success rate or increasing query cost.
The methodology therefore treats attack variants as comparable only when success, cost, and explanation behaviour are considered together.

\subsection{Dataset, model, and preprocessing}
All experiments use the CIFAR-10 test set \cite{krizhevsky2009cifar}, restricted to the fixed index range 0--9{,}999.
Using a fixed 10k index set keeps the evaluation stable across reruns and avoids ambiguity about which data slice the final report refers to.
The classifier is a pretrained CIFAR-10 ResNet-20 variant, used unchanged throughout the study.

Inputs are $32 \times 32$ RGB images, and the same preprocessing and normalisation pipeline is applied across all runs.
An image is considered \emph{eligible} if the clean model classifies it correctly.
This eligibility rule is important because the attack is intended to overturn an initially correct prediction rather than merely produce an arbitrary output on an already misclassified image.
All attack success rates reported later are therefore computed over the eligible subset rather than over the full 10k image range.

\subsection{Threat model}
The attacker is black-box: model parameters, gradients, and internal activations are unavailable.
The attack may interact with the classifier only through forward queries that return output scores.
The perturbation budget is restricted to one pixel location.
An adversarial example therefore differs from the clean image at exactly one coordinate $(x,y)$, although all three RGB values at that location may be changed.

This threat model is intentionally narrow.
Its purpose is not to emulate every realistic adversarial scenario, but to isolate sparse black-box optimisation in a setting where the perturbation is easy to inspect and the search space remains compact.
Within this model, the cost of an attack is measured in \emph{queries}, where one query corresponds to one forward evaluation of the classifier on a candidate perturbed image.

\subsection{Attack parameterisation}
All attack variants are based on the one-pixel formulation introduced by Su et al.\ \cite{su2019onepixel} and use Differential Evolution (DE) as the optimisation backbone \cite{storn1997de}.
A candidate solution encodes five continuous quantities: a pixel position $(x,y)$ and colour values $(r,g,b)$.
DE evolves a population of such candidates through mutation, crossover, and selection.

This parameterisation is suitable for the one-pixel setting for two reasons.
First, it is naturally derivative-free and therefore compatible with the black-box constraint.
Second, it turns the attack into a bounded low-dimensional search problem in which different search priors can be compared cleanly.
The report studies three attack configurations under this common representation.

\subsection{Untargeted one-pixel DE baseline}
The first configuration is an untargeted DE baseline.
Its objective is to induce any misclassification with a one-pixel change.
For the frozen 10k-scale run used in the report, the configuration uses population size 400, differential weight $F=0.5$, crossover rate $Cr=0.9$, maximum generations 100, early-stopping thresholds \texttt{es\_target}=0.9 and \texttt{es\_nonadv}=0.05, and random seed 1337.

This baseline plays two roles.
For RQ1, it defines the successful adversarial subset on which explanation changes are first analysed.
For RQ2, it provides the main unguided reference point against which the guided variants are compared.
The output records, for each image, the success flag, generation count, query count, and the best one-pixel parameters found.

\subsection{Fastprior multi-target DE}
The second configuration is a fixed-budget multi-target variant referred to as \emph{fastprior}.
Instead of attempting only an untargeted label change, the attack considers labels different from the clean prediction and runs DE with a fixed budget per target.
For the frozen 10k-scale run, the recorded configuration uses \texttt{pop}=192, maximum generations $=40$, $F=0.5$, $Cr=0.9$, \texttt{es\_target}$=0.9$, \texttt{xy\_topk}$=32$, \texttt{xy\_mutate\_prob}$=0.1$, and seed 1337.

This variant is useful because it changes the structure of the search itself.
Its near-constant query cost reflects the fixed-budget design, and its multi-target framing allows both image-wise and target-wise success to be measured.
The resulting trade-off is analytically interesting even if the configuration is not competitive with the untargeted baseline in headline efficiency.

\subsection{Responsibility-map-guided DE}
The third configuration is a responsibility-map-guided DE variant.
Here, the attack uses a black-box saliency prior derived from RISE \cite{petsiuk2018rise} to bias which pixel locations are explored during DE.
The purpose is not to change the one-pixel budget, but to change where the search spends effort.
If pixels that are influential for the current prediction are also useful places to perturb, then such guidance may improve the success--cost trade-off.

This is an empirical proposition rather than a guaranteed one.
A saliency map explains the current decision, not necessarily the cheapest route to changing that decision.
The guided variant therefore tests whether an intuitively appealing prior is actually useful under the constraints of sparse black-box search.

\subsection{Explanation methods}
To analyse explanation behaviour, three explanation methods are computed for paired clean and adversarial inputs:

\begin{itemize}
  \item \textbf{Grad-CAM} \cite{selvaraju2017gradcam}, which produces a class-specific localisation map using gradients over a late convolutional representation;
  \item \textbf{Integrated Gradients (IG)} \cite{sundararajan2017ig}, which computes attributions by integrating gradients from a baseline input to the observed input;
  \item \textbf{RISE} \cite{petsiuk2018rise}, which estimates black-box saliency by random masking and aggregation.
\end{itemize}

These methods are intentionally heterogeneous.
Grad-CAM is spatially coarse and architecture-dependent.
IG is gradient-based and path-dependent.
RISE is black-box, stochastic, and query-based.
Together they provide a useful spread of explanation paradigms for testing whether explanation changes under attack are method-specific or broadly shared.

\subsection{Metrics}
\subsubsection{Attack effectiveness}
The principal effectiveness metric is \emph{attack success rate} (ASR), computed over the eligible subset.
For fastprior, two effectiveness views are relevant: image-wise success, which asks whether at least one target succeeds for an image, and target-wise success, which counts successful targets across all target attempts.

\subsubsection{Attack cost}
Attack cost is measured in model queries.
The report emphasises median query counts over both all eligible images and successful images.
This is preferable to a mean-only summary because cost distributions in adversarial search are often skewed and can be heavily influenced by failures or long tails.

\subsubsection{Explanation stability}
Explanation stability is measured on successful adversarial examples using two complementary statistics:

\begin{itemize}
  \item \textbf{IoU@10}, the intersection-over-union of the top-$k$ salient pixels before and after attack, with $k=10$;
  \item \textbf{Spearman rank correlation} $\rho$, computed between flattened before/after attribution maps.
\end{itemize}

IoU@10 captures whether the most salient pixels remain similar, while Spearman $\rho$ captures whether the broader attribution ordering remains similar.
Using both helps reduce the risk that one metric alone overstates or understates stability.

\subsubsection{Explanation faithfulness}
Faithfulness is assessed using deletion AUC.
Pixels are removed in descending order of saliency, and the classifier score for the relevant class is tracked as more of the image is deleted.
The report presents clean and adversarial deletion AUC values together with the difference $\Delta$ deletion AUC (adversarial minus clean).
A more negative value indicates a larger reduction in the measured faithfulness of the explanation under attack.

\subsection{Aggregation protocol}
All runs use the same fixed 10k image range, and eligibility is determined once using the clean model.
Attack metrics are computed over the resulting eligible subset.
Explanation metrics are computed only on successful adversarial examples.
This means that explanation comparisons reflect both the perturbation mechanism and the composition of the successful subset produced by each attack.
That selection effect is not an accident of the analysis but part of the phenomenon being studied, although it must also be interpreted carefully later in \Cref{sec:threats_to_validity}.

\subsection{Frozen artefacts and report consistency}
To keep the report internally consistent, all final numbers are taken from the frozen results pack under \texttt{results/final\_pack/}.
The tables and plots used in the report are generated from stable CSV artefacts rather than copied manually.
This makes the report easier to audit and reduces the risk of transcription error.
The methodology therefore supports not only controlled comparison but also traceable reporting.
